% ============================================================
% DOCUMENT 3 : GUIDE TECHNIQUE D'ARCHITECTURE
% Titre : "Architecture Technique de l'Écosystème eSchool/MonEcole"
% Public : Ingénieurs, développeurs, contributeurs
% ============================================================
\documentclass[12pt,a4paper,oneside]{report}

% --- Encodage et langue ---
\usepackage[utf8]{inputenc}
\usepackage[T1]{fontenc}
\usepackage[french]{babel}
\usepackage{lmodern}

% --- Mise en page ---
\usepackage[margin=2.5cm]{geometry}
\usepackage{setspace}
\onehalfspacing
\usepackage{parskip}
\usepackage{fancyhdr}
\usepackage{enumitem}

% --- Graphiques ---
\usepackage{graphicx}
\usepackage[dvipsnames,svgnames,x11names]{xcolor}
\usepackage{tikz}
\usetikzlibrary{shapes.geometric, arrows.meta, positioning, calc, fit, backgrounds, shadows, decorations.pathmorphing, patterns, matrix, chains}

% --- Tableaux ---
\usepackage{booktabs}
\usepackage{tabularx}
\usepackage{longtable}
\usepackage{multirow}
\usepackage{array}

% --- Code et listings ---
\usepackage{listings}
\lstdefinestyle{python}{
  language=Python,
  basicstyle=\ttfamily\footnotesize,
  backgroundcolor=\color{gray!5},
  frame=single,
  rulecolor=\color{gray!30},
  breaklines=true,
  keywordstyle=\color{blue!70!black}\bfseries,
  stringstyle=\color{red!60!black},
  commentstyle=\color{green!50!black}\itshape,
  showstringspaces=false,
  tabsize=4,
  numbers=left,
  numberstyle=\tiny\color{gray},
  numbersep=5pt
}
\lstdefinestyle{sql}{
  language=SQL,
  basicstyle=\ttfamily\footnotesize,
  backgroundcolor=\color{blue!3},
  frame=single,
  rulecolor=\color{blue!20},
  breaklines=true,
  keywordstyle=\color{blue!70!black}\bfseries,
  stringstyle=\color{red!60!black},
  showstringspaces=false,
  morekeywords={VIEW, REPLACE, ENGINE, InnoDB, AUTO_INCREMENT, UNSIGNED}
}
\lstdefinestyle{js}{
  language=Java,
  basicstyle=\ttfamily\footnotesize,
  backgroundcolor=\color{yellow!3},
  frame=single,
  rulecolor=\color{yellow!30},
  breaklines=true,
  keywordstyle=\color{purple!70!black}\bfseries,
  stringstyle=\color{red!60!black},
  commentstyle=\color{green!50!black}\itshape,
  showstringspaces=false,
  morekeywords={async, await, const, let, fetch, function, document, getElementById, innerHTML}
}

% --- Boîtes tcbox ---
\usepackage[most]{tcolorbox}
\newtcolorbox{techbox}[1][]{enhanced, colback=blue!3, colframe=blue!50!black, fonttitle=\bfseries, title=#1}
\newtcolorbox{warnbox}[1][]{enhanced, colback=orange!5, colframe=orange!60!black, fonttitle=\bfseries, title=#1}
\newtcolorbox{patternbox}[1][]{enhanced, colback=green!3, colframe=green!50!black, fonttitle=\bfseries, title=#1}
\newtcolorbox{pitfallbox}[1][]{enhanced, colback=red!3, colframe=red!50!black, fonttitle=\bfseries, title=#1}

% --- Hyperliens ---
\usepackage[colorlinks=true,linkcolor=blue!60!black,urlcolor=blue!50!black]{hyperref}
\usepackage{bookmark}

% --- En-têtes ---
\pagestyle{fancy}
\fancyhf{}
\fancyhead[L]{\small\leftmark}
\fancyhead[R]{\small Architecture Technique eSchool}
\fancyfoot[C]{\thepage}
\renewcommand{\headrulewidth}{0.4pt}

% --- Couleurs ---
\definecolor{eschoolblue}{HTML}{2563EB}
\definecolor{eschoolpurple}{HTML}{7C3AED}
\definecolor{eschoolteal}{HTML}{0D9488}
\definecolor{eschoolamber}{HTML}{D97706}
\definecolor{hubcolor}{HTML}{3B82F6}
\definecolor{spokecolor}{HTML}{10B981}
\definecolor{arrowcolor}{HTML}{6366F1}
\definecolor{viewcolor}{HTML}{F59E0B}

% ============================================================
\begin{document}

% --- PAGE DE COUVERTURE ---
\begin{titlepage}
\begin{tikzpicture}[remember picture, overlay]
  \fill[left color=gray!90!black, right color=eschoolblue!80!black] 
    (current page.south west) rectangle (current page.north east);
  \foreach \x/\y/\r/\o in {3/24/3/0.04, 14/22/4/0.05, 6/6/5/0.03, 17/4/2.5/0.06} {
    \fill[white, opacity=\o] (\x,\y) circle (\r);
  }
  \fill[white, opacity=0.95] 
    ([yshift=5cm]current page.south west) rectangle (current page.south east);
\end{tikzpicture}

\vspace*{3cm}
\begin{center}
  {\Huge\bfseries\color{white} Architecture Technique}\\[0.5cm]
  {\Huge\bfseries\color{white} de l'Écosystème}\\[0.5cm]
  {\Huge\bfseries\color{eschoolblue!30!white} eSchool / MonEcole}\\[1cm]
  {\Large\color{white!80} Guide Technique Approfondi}\\[0.3cm]
  {\large\color{white!70} Pour Ingénieurs et Développeurs}\\[4cm]
  \rule{10cm}{0.5pt}\\[0.5cm]
  {\color{white!80} Django $\bullet$ MySQL/MariaDB $\bullet$ ReportLab $\bullet$ Gunicorn}\\[5cm]
\end{center}
\begin{center}
  {\normalsize\color{gray!80} Version 1.0 — Mars 2026}\\[0.2cm]
  {\normalsize\color{gray!80} evanet08 — \texttt{github.com/evanet08}}
\end{center}
\end{titlepage}

\tableofcontents
\newpage

% ============================================================
\chapter{Vue d'Ensemble Technique}
% ============================================================

\section{Stack Technologique}

\begin{center}
\begin{tikzpicture}[
  layer/.style={rectangle, rounded corners=6pt, draw, minimum width=13cm, minimum height=1.8cm, align=center, font=\small, drop shadow={shadow xshift=0.5pt, shadow yshift=-0.5pt}},
  arr/.style={-{Stealth[length=5pt]}, thick, gray!40}
]
  \node[layer, fill=eschoolblue!10, draw=eschoolblue!40] (client) at (0,7) {
    \textbf{Client} — Navigateur Web\\
    HTML5 $\bullet$ CSS3 (Variables, Glassmorphism) $\bullet$ JavaScript (Vanilla, AJAX/Fetch)\\
    Bootstrap 5 $\bullet$ FontAwesome 6 $\bullet$ Google Fonts (Inter)
  };
  
  \node[layer, fill=eschoolpurple!10, draw=eschoolpurple!40] (web) at (0,4.5) {
    \textbf{Serveur Web} — Apache 2.4 + mod\_proxy\\
    Reverse proxy vers Gunicorn $\bullet$ SSL/TLS via Plesk\\
    Static files servis directement par Apache
  };
  
  \node[layer, fill=eschoolteal!10, draw=eschoolteal!40] (app) at (0,2) {
    \textbf{Application} — Django 4.2 + Gunicorn\\
    2 projets : eSchoolStructure (Hub) + MonEcole (Spoke)\\
    ReportLab (PDF) $\bullet$ openpyxl (Excel) $\bullet$ Brevo (Email)
  };
  
  \node[layer, fill=eschoolamber!10, draw=eschoolamber!40] (db) at (0,-0.5) {
    \textbf{Base de Données} — MySQL 8.0 / MariaDB 10.6\\
    \texttt{countryStructure} (Hub) + \texttt{db\_monecole} (Spoke)\\
    Cross-DB VIEWs $\bullet$ DatabaseRouter $\bullet$ \texttt{managed=False}
  };
  
  \draw[arr] (client) -- (web);
  \draw[arr] (web) -- (app);
  \draw[arr] (app) -- (db);
\end{tikzpicture}
\end{center}

\section{Repositories Git}

\begin{center}
\begin{tabular}{l l l}
\toprule
\textbf{Repository} & \textbf{Rôle} & \textbf{Serveur} \\
\midrule
\texttt{evanet08/eSchoolStructure} & Hub national & \texttt{eschool-rdc.pro} \\
\texttt{evanet08/MonEcole} & Spoke scolaire & \texttt{monecole.pro} \\
\bottomrule
\end{tabular}
\end{center}

\section{Arborescence des Projets}

\begin{lstlisting}[style=python, title=eSchoolStructure]
eSchoolStructure/
|-- eSchoolStructure/       # Projet Django principal
|   |-- settings.py         # Config BDD, SMTP, static
|   |-- urls.py             # Routage principal
|   |-- db_routers.py       # DatabaseRouter
|-- structure_app/          # Application principale
|   |-- models.py           # 25+ modeles (1145 lignes)
|   |-- views.py            # 80+ endpoints (2600+ lignes)
|   |-- auth_views.py       # Auth hierarchique
|   |-- templates/          # Templates Django
|   |-- static/             # CSS, JS, images
|   |-- migrations/         # 28 migrations
|-- docs/latex/             # Cette documentation
|-- manage.py
\end{lstlisting}

\begin{lstlisting}[style=python, title=MonEcole]
MonEcole/
|-- MonEcole/               # Projet Django principal
|   |-- settings.py
|   |-- urls.py
|   |-- db_routers.py       # CountryStructureRouter
|-- MonEcole_app/
|   |-- models/             # Modeles repartis en modules
|   |   |-- country_structure.py  # Tables Hub (managed=False)
|   |   |-- classe.py       # Classes, cycles, sections
|   |   |-- eleves/         # Eleve, Inscription, Notes
|   |   |-- evaluations/    # Notes, Deliberations
|   |   |-- recouvrement/   # Paiements, Rubriques
|   |-- views/              # Views reparties en modules
|   |   |-- direction_views/ # Dashboard directeur
|   |   |-- inscription/    # CRUD eleves
|   |   |-- evaluation/     # Notes et deliberations
|   |   |-- enseignement/   # Cours et attributions
|   |   |-- recouvrement/   # Finances
|   |-- templates/          # 60+ templates HTML
\end{lstlisting}


% ============================================================
\chapter{Architecture Multi-Base de Données}
% ============================================================

\section{Vue d'Ensemble du Schéma Distribué}

\begin{center}
\begin{tikzpicture}[
  db/.style={cylinder, draw, shape border rotate=90, minimum width=2.5cm, minimum height=2cm, font=\small\bfseries, align=center, aspect=0.25, drop shadow},
  app/.style={rectangle, rounded corners=6pt, draw, minimum width=3.5cm, minimum height=1.5cm, font=\small, align=center, drop shadow},
  view/.style={rectangle, rounded corners=3pt, draw=viewcolor!60, fill=viewcolor!10, minimum width=2cm, minimum height=0.5cm, font=\tiny, align=center, dashed},
  arr/.style={-{Stealth[length=5pt]}, thick},
  darr/.style={-{Stealth[length=5pt]}, thick, dashed}
]
  % Hub
  \node[db, fill=hubcolor!15, draw=hubcolor!60] (dbhub) at (0,0) {\texttt{country}\\\texttt{Structure}};
  \node[app, fill=hubcolor!8, draw=hubcolor!50] (apphub) at (0,3) {eSchool\\Structure};
  
  % Spoke
  \node[db, fill=spokecolor!15, draw=spokecolor!60] (dbspoke) at (8,0) {\texttt{db\_}\\\texttt{monecole}};
  \node[app, fill=spokecolor!8, draw=spokecolor!50] (appspoke) at (8,3) {MonEcole};
  
  % VIEWs
  \node[view] (v1) at (4,1.2) {VIEW cours};
  \node[view] (v2) at (4,0.4) {VIEW trimestres};
  \node[view] (v3) at (4,-0.4) {VIEW classes};
  \node[view] (v4) at (4,-1.2) {VIEW ecole};
  
  % Connexions
  \draw[arr, hubcolor] (apphub) -- (dbhub) node[midway, left, font=\scriptsize] {ORM};
  \draw[arr, spokecolor] (appspoke) -- (dbspoke) node[midway, right, font=\scriptsize] {ORM};
  
  \draw[darr, viewcolor!70] (dbhub) -- (v1);
  \draw[darr, viewcolor!70] (dbhub) -- (v2);
  \draw[darr, viewcolor!70] (dbhub) -- (v3);
  \draw[darr, viewcolor!70] (dbhub) -- (v4);
  \draw[darr, viewcolor!70] (v1) -- (dbspoke);
  \draw[darr, viewcolor!70] (v2) -- (dbspoke);
  \draw[darr, viewcolor!70] (v3) -- (dbspoke);
  \draw[darr, viewcolor!70] (v4) -- (dbspoke);
  
  % Légende
  \node[font=\scriptsize, text=viewcolor!80!black] at (4,-2.2) {VIEWs MySQL cross-BDD};
\end{tikzpicture}
\end{center}

\section{Le DatabaseRouter}

Le \texttt{CountryStructureRouter} aiguille les requêtes ORM vers la bonne base :

\begin{lstlisting}[style=python, title=db\_routers.py]
class CountryStructureRouter:
    """
    Route les modeles managed=False vers countryStructure,
    tout le reste vers la base par defaut (db_monecole).
    """
    hub_models = {
        'session', 'trimestre', 'periode', 'mention',
        'pays', 'classe', 'cycle',
    }
    
    def db_for_read(self, model, **hints):
        if model._meta.db_table in self.hub_models:
            return 'countryStructure'
        if not model._meta.managed:
            return 'countryStructure'
        return 'default'
    
    def db_for_write(self, model, **hints):
        if not model._meta.managed:
            return 'countryStructure'
        return 'default'
    
    def allow_relation(self, obj1, obj2, **hints):
        return True  # Cross-DB relations autorisees
    
    def allow_migrate(self, db, app_label, model_name=None, **hints):
        return db == 'default'
\end{lstlisting}

\section{Pattern managed=False}

Les tables du Hub réutilisées dans le Spoke sont déclarées avec \texttt{managed = False} :

\begin{lstlisting}[style=python, title=MonEcole models/country\_structure.py]
class Trimestre(models.Model):
    """
    Table de reference : countryStructure.trimestres
    managed = False : Django ne modifie pas le schema.
    """
    id_trimestre = models.AutoField(primary_key=True)
    trimestre = models.CharField(max_length=50, unique=True)
    is_active = models.BooleanField(default=True)
    
    class Meta:
        db_table = 'trimestres'
        managed = False  # <-- Cle du pattern
        verbose_name = 'Trimestre'
    
    def __str__(self):
        return self.trimestre
\end{lstlisting}

\section{VIEWs MySQL}

\subsection{VIEWs de Référence Simples}

\begin{lstlisting}[style=sql, title=VIEWs 1:1 vers le Hub]
-- Trimestres
CREATE OR REPLACE VIEW `db_monecole`.`trimestres` AS
SELECT * FROM countryStructure.trimestres;

-- Sessions
CREATE OR REPLACE VIEW `db_monecole`.`sessions` AS
SELECT * FROM countryStructure.sessions;

-- Mentions
CREATE OR REPLACE VIEW `db_monecole`.`mentions` AS
SELECT * FROM countryStructure.mentions;

-- Classes
CREATE OR REPLACE VIEW `db_monecole`.`classes` AS
SELECT * FROM countryStructure.classes;
\end{lstlisting}

\subsection{VIEW Cross-DB avec JOIN : cours\_par\_classe}

\begin{lstlisting}[style=sql, title=VIEW avec jointure cross-BDD]
CREATE OR REPLACE VIEW `db_monecole`.`cours_par_classe` AS
SELECT 
    c.id_cours AS id_cours_classe,
    c.cours,
    c.code_cours,
    c.classe_id,
    c.domaine_id,
    ca.id_classe_active,
    ca.id_campus_id,
    ca.id_annee_id
FROM countryStructure.cours c
JOIN db_monecole.classe_active ca 
  ON ca.classe_id_id = c.classe_id;
\end{lstlisting}

\subsection{VIEW Universal Context : annee\_trimestre}

\begin{pitfallbox}[Piège : Le Pattern ``Hardcoded Zero'']
Ne JAMAIS utiliser \texttt{0 AS id\_campus\_id} dans une VIEW. Cela casse tous les filtres Django ORM. Toujours effectuer un JOIN pour résoudre les identifiants de contexte.
\end{pitfallbox}

\begin{lstlisting}[style=sql, title=Pattern ``Universal Context'']
CREATE OR REPLACE VIEW `db_monecole`.`annee_trimestre` AS
SELECT 
    t.id AS id_trimestre,
    t.trimestre_id,
    t.debut, t.fin,
    t.etat AS etat_trimestre,
    a.id_annee AS id_annee_id,
    c.id_campus AS id_campus_id,
    cca.id_cycle_actif AS id_cycle_id,
    ca.id_classe_active AS id_classe_id
FROM countryStructure.etablissements_annees_trimestres t
JOIN countryStructure.etablissements_annees ea 
  ON t.etablissement_annee_id = ea.id
JOIN db_monecole.annee a 
  ON a.id_annee_structure = ea.annee_id
JOIN db_monecole.campus c 
  ON c.id_etablissement = ea.etablissement_id 
  AND c.is_active = 1
JOIN db_monecole.classe_cycle_actif cca 
  ON cca.id_campus_id = c.id_campus 
  AND cca.id_annee_id = a.id_annee 
  AND cca.is_active = 1
JOIN db_monecole.classe_active ca 
  ON ca.cycle_id_id = cca.id_cycle_actif 
  AND ca.id_campus_id = c.id_campus 
  AND ca.id_annee_id = a.id_annee 
  AND ca.is_active = 1;
\end{lstlisting}


% ============================================================
\chapter{Schéma de Données — eSchoolStructure}
% ============================================================

\section{Diagramme Entité-Relation Complet}

\begin{center}
\begin{tikzpicture}[
  entity/.style={rectangle, rounded corners=3pt, draw=#1!60, fill=#1!8, minimum width=3cm, minimum height=0.7cm, font=\scriptsize\bfseries, align=center},
  fk/.style={-{Stealth[length=3pt]}, thin, gray!60}
]
  % === COUCHE 1 : PAYS ===
  \node[entity=red] (pays) at (0,0) {Pays};
  
  % === COUCHE 2 : SOUS-PAYS ===
  \node[entity=orange] (pedtype) at (-6,-1.5) {PedagogicStructureType};
  \node[entity=orange] (admtype) at (-2.5,-1.5) {AdministrativeStructureType};
  \node[entity=eschoolblue] (cycle) at (1,-1.5) {Cycle};
  \node[entity=eschoolpurple] (regime) at (4,-1.5) {Regime};
  \node[entity=eschoolteal] (programme) at (7,-1.5) {Programme};
  
  % === COUCHE 3 : INSTANCES ===
  \node[entity=orange] (pedinst) at (-6,-3) {PedagogicStructureInstance};
  \node[entity=orange] (adminst) at (-2.5,-3) {AdminStructureInstance};
  \node[entity=eschoolblue] (section) at (1,-3) {Section};
  \node[entity=eschoolblue] (classe) at (4,-3) {Classe};
  \node[entity=eschoolteal] (domaine) at (7,-3) {Domaine};
  
  % === COUCHE 4 : COURS ===
  \node[entity=hubcolor] (cours) at (2,-4.5) {Cours};
  \node[entity=hubcolor] (coursannee) at (6,-4.5) {CoursAnnee};
  
  % === COUCHE 5 : ÉTABLISSEMENTS ===
  \node[entity=spokecolor] (gest) at (-4,-4.5) {GestionnaireEtab};
  \node[entity=spokecolor] (etab) at (-1,-6) {Etablissement};
  \node[entity=eschoolamber] (annee) at (3,-6) {Annee};
  
  % === COUCHE 6 : CONFIG ANNUELLE ===
  \node[entity=eschoolamber] (etabannee) at (0,-7.5) {EtablissementAnnee};
  \node[entity=eschoolamber] (etabclasse) at (5,-7.5) {EtabAnneeClasse};
  
  % === COUCHE 7 : TEMPORAL ===
  \node[entity=arrowcolor] (session) at (-5,-7.5) {Session};
  \node[entity=arrowcolor] (trimestre) at (-3,-9) {Trimestre};
  \node[entity=arrowcolor] (periode) at (0,-9) {Periode};
  \node[entity=arrowcolor] (mention) at (3,-9) {Mention};
  \node[entity=arrowcolor] (etabtrim) at (-1,-10.5) {EtabAnneeTrimestre};
  \node[entity=arrowcolor] (etabper) at (4,-10.5) {EtabAnneePeriode};
  
  % === COUCHE 8 : SECURITE ===
  \node[entity=red] (admin) at (-5,-6) {AdminUser};
  \node[entity=red] (otp) at (-7,-7.5) {OTPCode};
  
  % === FK ARROWS ===
  \foreach \child in {pedtype,admtype,cycle,regime,programme} {\draw[fk](\child)--(pays);}
  \draw[fk] (classe) -- (cycle);
  \draw[fk] (cours) -- (classe);
  \draw[fk] (cours) -- (domaine);
  \draw[fk] (cours) -- (section);
  \draw[fk] (coursannee) -- (cours);
  \draw[fk] (coursannee) -- (annee);
  \draw[fk] (coursannee) -- (domaine);
  \draw[fk] (etab) -- (pays);
  \draw[fk] (etab) -- (gest);
  \draw[fk] (etabannee) -- (etab);
  \draw[fk] (etabannee) -- (annee);
  \draw[fk] (etabclasse) -- (etabannee);
  \draw[fk] (etabclasse) -- (classe);
  \draw[fk] (etabclasse) -- (section);
  \draw[fk] (etabtrim) -- (etabannee);
  \draw[fk] (etabtrim) -- (trimestre);
  \draw[fk] (etabper) -- (etabtrim);
  \draw[fk] (etabper) -- (periode);
  \draw[fk] (admin) -- (pays);
  \draw[fk] (admin) -- (etab);
  \draw[fk] (admin) -- (pedinst);
  \draw[fk] (otp) -- (admin);
  \draw[fk] (periode) -- (trimestre);
\end{tikzpicture}
\end{center}

\section{Modèles Détaillés}

\subsection{Pays — Point d'Entrée}

\begin{center}
\begin{tabularx}{\textwidth}{l l X}
\toprule
\textbf{Champ} & \textbf{Type} & \textbf{Description} \\
\midrule
\texttt{id\_pays} & AutoField (PK) & Identifiant unique \\
\texttt{nom} & CharField(100) & Nom du pays (unique) \\
\texttt{sigle} & CharField(5) & Code ISO du pays (unique) \\
\texttt{nLevelsStructuraux} & PositiveInt & Nombre de niveaux pédagogiques \\
\texttt{nLevelsAdministratifs} & PositiveInt & Nombre de niveaux administratifs \\
\texttt{domaine} & CharField & Domaine web associé \\
\bottomrule
\end{tabularx}
\end{center}

\subsection{Cycle — Niveaux d'Enseignement}

\begin{center}
\begin{tabularx}{\textwidth}{l l X}
\toprule
\textbf{Champ} & \textbf{Type} & \textbf{Description} \\
\midrule
\texttt{id\_cycle} & AutoField (PK) & Identifiant unique \\
\texttt{nom} & CharField(100) & Nom (Primaire, Secondaire...) \\
\texttt{code} & CharField(20) & Code court (PRI, SEC...) \\
\texttt{duree} & PositiveInt & Durée en années \\
\texttt{ordre} & PositiveInt & Ordre hiérarchique \\
\texttt{pays} & FK → Pays & Pays d'appartenance \\
\texttt{hasSections} & Boolean & Possède des sections (Sci, Lit...) \\
\bottomrule
\end{tabularx}
\end{center}

\subsection{Cours — Catalogue des Matières}

\begin{center}
\begin{tabularx}{\textwidth}{l l X}
\toprule
\textbf{Champ} & \textbf{Type} & \textbf{Description} \\
\midrule
\texttt{id\_cours} & AutoField (PK) & Identifiant unique \\
\texttt{cours} & CharField(200) & Nom du cours (Mathématiques...) \\
\texttt{code\_cours} & CharField(20) & Code unique (MATH01...) \\
\texttt{classe} & FK → Classe & Classe d'appartenance \\
\texttt{domaine} & FK → Domaine & Domaine par défaut \\
\texttt{section} & FK → Section (null) & Section (pour cycles à sections) \\
\multicolumn{3}{l}{\textit{unique\_together: (classe, section, code\_cours)}} \\
\bottomrule
\end{tabularx}
\end{center}

\subsection{AdminUser — Hiérarchie de Sécurité}

\begin{center}
\begin{tabularx}{\textwidth}{l l X}
\toprule
\textbf{Champ} & \textbf{Type} & \textbf{Description} \\
\midrule
\texttt{id\_admin} & AutoField (PK) & Identifiant unique \\
\texttt{email} & EmailField (unique) & Identifiant principal \\
\texttt{telephone} & CharField(20) & Numéro international \\
\texttt{password\_hash} & CharField(128) & Hash PBKDF2-SHA256 \\
\texttt{pays} & FK → Pays & Pays d'appartenance \\
\texttt{structure\_instance} & FK → PedagogicInst (null) & Niveau structurel \\
\texttt{etablissement} & FK → Etablissement (null) & École spécifique \\
\texttt{email\_verified} & Boolean & Email vérifié par OTP \\
\texttt{is\_active} & Boolean & Compte actif \\
\texttt{created\_by} & FK → self (null) & Administrateur créateur \\
\bottomrule
\end{tabularx}
\end{center}

\begin{techbox}[Propriétés Dynamiques de AdminUser]
\begin{itemize}[topsep=0pt, itemsep=2pt]
  \item \texttt{niveau\_ordre} — 0 (National), 1-N (Structure), N+1 (École)
  \item \texttt{niveau\_nom} — ``National'', ``Niveau X'', ``Établissement''
  \item \texttt{scope\_code} — Code hiérarchique pour filtrage par préfixe
  \item \texttt{scope\_name} — Nom de l'entité gérée
  \item \texttt{is\_national} — True si niveau 0
  \item \texttt{is\_validated} — True si email ou phone vérifié
  \item \texttt{has\_password} — True si mot de passe défini
\end{itemize}
\end{techbox}


% ============================================================
\chapter{Schéma de Données — MonEcole}
% ============================================================

\section{Vue d'Ensemble}

MonEcole utilise 30+ modèles répartis dans des modules Python :

\begin{center}
\begin{tikzpicture}[
  mod/.style={rectangle, rounded corners=4pt, draw=#1!50, fill=#1!8, minimum width=4.5cm, minimum height=1.8cm, font=\scriptsize, align=left},
]
  \node[mod=spokecolor] (m1) at (0,3) {\textbf{country\_structure.py}\\Pays, Session, Trimestre,\\Periode, Mention (managed=False)};
  \node[mod=eschoolblue] (m2) at (6,3) {\textbf{classe.py}\\Classe, Cycle, Classe\_active,\\Section, Délibération};
  \node[mod=eschoolpurple] (m3) at (0,0) {\textbf{eleves/}\\Eleve, Inscription, Note\_type,\\Eleve\_note, Conduite};
  \node[mod=eschoolamber] (m4) at (6,0) {\textbf{evaluations/}\\Deliberation\_type, Finalite,\\5 types de résultats, Evaluation};
  \node[mod=red] (m5) at (0,-3) {\textbf{recouvrement/}\\Paiement, Rubrique\_variable,\\Prix, Derogation};
  \node[mod=eschoolteal] (m6) at (6,-3) {\textbf{Autres}\\Campus, Ecole, Annee, Module,\\Personnel, Horaire, Salle};
\end{tikzpicture}
\end{center}

\section{Modèle Élève — Fiche Complète}

Le modèle \texttt{Eleve} contient 35+ champs couvrant l'identité, la filiation, la localisation et les identifiants de sécurité :

\begin{center}
\begin{tabularx}{\textwidth}{l X}
\toprule
\textbf{Catégorie} & \textbf{Champs} \\
\midrule
Identité & nom, prenom, genre, etat\_civil, date\_naissance \\
Filiation & nom\_pere, prenom\_pere, nom\_mere, prenom\_mere, tutaire \\
Scolarité & code\_eleve, code\_annee, matricule \\
Contact & email, email\_parent, telephone \\
Localisation & naissance\_pays/province/commune/zone, province/commune/zone\_actuelle \\
Sécurité & password, password\_parent, code\_secret\_parent, code\_secret\_eleve \\
Media & imageUrl (upload vers logos/eleves/) \\
\bottomrule
\end{tabularx}
\end{center}

\section{Circuit de la Note}

\begin{center}
\begin{tikzpicture}[
  step/.style={rectangle, rounded corners=5pt, draw=#1!50, fill=#1!10, minimum width=2.8cm, minimum height=1.2cm, font=\scriptsize, align=center},
  arr/.style={-{Stealth[length=5pt]}, thick, arrowcolor!60}
]
  \node[step=eschoolblue] (s1) at (0,0) {\textbf{Enseignant}\\Saisit les notes};
  \node[step=eschoolpurple] (s2) at (4,0) {\textbf{Evaluation}\\Type + Période\\+ Session};
  \node[step=eschoolteal] (s3) at (8,0) {\textbf{Eleve\_note}\\Note/20, Note/Max\\Repêchage};
  \node[step=eschoolamber] (s4) at (12,0) {\textbf{Délibération}\\Pourcentage +\\Rang + Mention};
  \node[step=red] (s5) at (8,-2.5) {\textbf{Bulletin PDF}\\ReportLab};
  
  \draw[arr] (s1) -- (s2);
  \draw[arr] (s2) -- (s3);
  \draw[arr] (s3) -- (s4);
  \draw[arr] (s4) -- (s5);
  \draw[arr] (s3) -- (s5);
\end{tikzpicture}
\end{center}


% ============================================================
\chapter{Liaison Hub $\leftrightarrow$ Spoke}
% ============================================================

\section{Chaîne de Résolution du Tenant}

\begin{center}
\begin{tikzpicture}[
  box/.style={rectangle, rounded corners=5pt, draw=#1!50, fill=#1!10, minimum width=3.2cm, minimum height=1.5cm, font=\scriptsize, align=center},
  arr/.style={-{Stealth[length=5pt]}, thick}
]
  \node[box=spokecolor] (campus) at (0,0) {\textbf{campus}\\(\texttt{db\_monecole})\\[2pt]\texttt{id\_etablissement}\\= lien vers Hub};
  
  \node[box=hubcolor] (etab) at (4.5,0) {\textbf{etablissements}\\(\texttt{countryStructure})\\[2pt]Données de l'école\\Matricule, Logo};
  
  \node[box=hubcolor] (ea) at (9,0) {\textbf{etab\_annees}\\(\texttt{countryStructure})\\[2pt]Année activée\\pour cette école};
  
  \node[box=arrowcolor] (trim) at (4.5,-3) {\textbf{etab\_trim}\\Trimestres\\configurés};
  
  \node[box=arrowcolor] (per) at (9,-3) {\textbf{etab\_per}\\Périodes\\configurées};
  
  \draw[arr, spokecolor!70] (campus) -- (etab) node[midway, above, font=\tiny\itshape] {\texttt{id\_etablissement}};
  \draw[arr, hubcolor!70] (etab) -- (ea) node[midway, above, font=\tiny\itshape] {FK};
  \draw[arr, arrowcolor!70] (ea) -- (trim);
  \draw[arr, arrowcolor!70] (trim) -- (per);
\end{tikzpicture}
\end{center}

\begin{techbox}[Règle d'Or : ``Always Carry the Etablissement ID'']
Chaque enregistrement opérationnel dans le Spoke doit avoir un chemin de résolution vers \texttt{id\_etablissement} via la table \texttt{campus}. Sans ce lien, les VIEWs cross-BDD retournent des ensembles vides.
\end{techbox}


% ============================================================
\chapter{Système d'Authentification}
% ============================================================

\section{Machine à États}

\begin{center}
\begin{tikzpicture}[
  state/.style={ellipse, draw=#1!60, fill=#1!10, minimum width=2.5cm, minimum height=1cm, font=\scriptsize\bfseries, align=center},
  arr/.style={-{Stealth[length=5pt]}, thick, gray!70},
  lab/.style={font=\tiny\itshape, text=gray!70}
]
  \node[state=eschoolblue] (s1) at (0,0) {Email\\Input};
  \node[state=eschoolpurple] (s2) at (5,1.5) {Password\\Entry};
  \node[state=eschoolamber] (s3) at (5,-1.5) {OTP\\Choice};
  \node[state=red] (s4) at (10,-1.5) {OTP\\Verify};
  \node[state=eschoolteal] (s5) at (10,1.5) {Set\\Password};
  \node[state=spokecolor] (s6) at (14,0) {Dashboard\\(Logged In)};
  
  \draw[arr] (s1) -- (s2) node[midway, lab, above] {verified + has\_pw};
  \draw[arr] (s1) -- (s3) node[midway, lab, below] {unverified};
  \draw[arr] (s3) -- (s4) node[midway, lab, below] {Email/SMS};
  \draw[arr] (s4) -- (s5) node[midway, lab, right] {code OK};
  \draw[arr] (s2) -- (s6) node[midway, lab, above] {pw OK};
  \draw[arr] (s5) -- (s6) node[midway, lab, above] {pw saved};
\end{tikzpicture}
\end{center}

\section{Endpoints d'Authentification}

\begin{center}
\begin{tabularx}{\textwidth}{l l X}
\toprule
\textbf{Endpoint} & \textbf{Méthode} & \textbf{Description} \\
\midrule
\texttt{/api/auth/check-email/} & POST & Vérifie existence et statut \\
\texttt{/api/auth/request-otp/} & POST & Génère et envoie OTP \\
\texttt{/api/auth/verify-otp/} & POST & Valide le code OTP \\
\texttt{/api/auth/set-password/} & POST & Définit le mot de passe \\
\texttt{/api/auth/login/} & POST & Connexion classique \\
\texttt{/api/auth/me/} & GET & Infos utilisateur connecté \\
\texttt{/api/auth/logout/} & POST & Déconnexion \\
\bottomrule
\end{tabularx}
\end{center}

\section{Décorateurs de Sécurité}

\begin{lstlisting}[style=python, title=Décorateurs @require\_auth et @require\_level]
def require_auth(view_func):
    """Verifie que la requete provient d'un AdminUser connecte."""
    def wrapper(request, *args, **kwargs):
        user_id = request.session.get('user_id')
        if not user_id:
            return JsonResponse({'error': 'Non authentifie'}, status=401)
        try:
            request.admin_user = AdminUser.objects.get(id_admin=user_id)
        except AdminUser.DoesNotExist:
            request.session.flush()
            return JsonResponse({'error': 'Session invalide'}, status=401)
        return view_func(request, *args, **kwargs)
    return wrapper

def require_level(max_level):
    """Verifie le niveau hierarchique minimum."""
    def decorator(view_func):
        def wrapper(request, *args, **kwargs):
            if request.admin_user.niveau_ordre > max_level:
                return JsonResponse({'error': 'Acces refuse'}, status=403)
            return view_func(request, *args, **kwargs)
        return wrapper
    return decorator
\end{lstlisting}


% ============================================================
\chapter{Patterns de Développement}
% ============================================================

\section{Pattern Upsert JSON}

\begin{patternbox}[Pattern : Upsert (Create or Update)]
Un seul endpoint gère à la fois la création et la modification d'une entité. La présence ou l'absence de l'ID dans le payload détermine l'opération.
\end{patternbox}

\begin{lstlisting}[style=python, title=Exemple : save\_cours]
@require_http_methods(["POST"])
def save_cours(request):
    data = json.loads(request.body)
    id_cours = data.get('id_cours')  # None = create
    
    payload = {
        'cours': data.get('cours'),
        'code_cours': data.get('code_cours'),
        'classe': get_object_or_404(Classe, id_classe=data['id_classe']),
        'domaine': get_object_or_404(Domaine, id_domaine=data['id_domaine']),
    }
    
    if id_cours:
        # UPDATE
        instance = get_object_or_404(Cours, id_cours=id_cours)
        for key, val in payload.items():
            setattr(instance, key, val)
        instance.save()
    else:
        # CREATE
        if Cours.objects.filter(
            classe=payload['classe'], 
            code_cours=payload['code_cours']
        ).exists():
            return JsonResponse({'error': 'Code duplique'}, status=400)
        instance = Cours.objects.create(**payload)
    
    return JsonResponse({'success': True, 'id': instance.id_cours})
\end{lstlisting}

\section{Pattern Checkbox Activation}

\begin{patternbox}[Pattern : Activation par Checkbox]
L'utilisateur voit TOUS les éléments du catalogue. Les éléments actifs dans le contexte courant sont cochés. Cocher/décocher déclenche une insertion/suppression atomique via AJAX.
\end{patternbox}

\begin{center}
\begin{tikzpicture}[
  item/.style={rectangle, draw=gray!30, fill=gray!5, minimum width=8cm, minimum height=0.8cm, font=\small, align=left},
]
  \node[item, fill=spokecolor!10] (i1) at (0,2) {~$\boxtimes$~~MATH01 — Mathématiques~~\textit{(Pondération: 3, Heures: 6h)}};
  \node[item, fill=spokecolor!10] (i2) at (0,1) {~$\boxtimes$~~FRAN01 — Français~~\textit{(Pondération: 2, Heures: 4h)}};
  \node[item] (i3) at (0,0) {~$\square$~~ANGL01 — Anglais~~\textit{(Non activé cette année)}};
  \node[item, fill=spokecolor!10] (i4) at (0,-1) {~$\boxtimes$~~PHYS01 — Physique~~\textit{(Pondération: 2, Heures: 3h)}};
  \node[item] (i5) at (0,-2) {~$\square$~~CHIM01 — Chimie~~\textit{(Non activé cette année)}};
  
  \node[font=\scriptsize\itshape, text=gray] at (0,-3) {Cocher = INSERT dans cours\_annee | Décocher = DELETE (si pas de notes liées)};
\end{tikzpicture}
\end{center}


% ============================================================
\chapter{Déploiement et DevOps}
% ============================================================

\section{Workflow de Déploiement}

\begin{center}
\begin{tikzpicture}[
  step/.style={rectangle, rounded corners=5pt, draw=eschoolblue!50, fill=eschoolblue!8, minimum width=3cm, minimum height=1.2cm, font=\scriptsize, align=center},
  arr/.style={-{Stealth[length=5pt]}, thick, eschoolblue!60}
]
  \node[step] (s1) at (0,0) {\textbf{1. Local}\\git add + commit\\+ push origin main};
  \node[step] (s2) at (4.5,0) {\textbf{2. SSH}\\Connexion au\\serveur production};
  \node[step] (s3) at (9,0) {\textbf{3. Pull}\\git fetch origin\\git reset --hard};
  \node[step] (s4) at (0,-2.5) {\textbf{4. Static}\\collectstatic\\--noinput};
  \node[step] (s5) at (4.5,-2.5) {\textbf{5. Migrate}\\manage.py\\migrate};
  \node[step] (s6) at (9,-2.5) {\textbf{6. Restart}\\systemctl restart\\eschool-structure};
  
  \draw[arr] (s1) -- (s2);
  \draw[arr] (s2) -- (s3);
  \draw[arr] (s3) -- (s4);
  \draw[arr] (s4) -- (s5);
  \draw[arr] (s5) -- (s6);
\end{tikzpicture}
\end{center}

\begin{lstlisting}[style=python, title=Commande de déploiement complète]
sshpass -p 'PASSWORD' ssh root@87.106.23.108 \
  "cd /var/www/vhosts/eschool-rdc.pro/httpdocs && \
   git fetch origin && \
   git reset --hard origin/main && \
   ./venv/bin/python3 manage.py collectstatic --noinput && \
   ./venv/bin/python3 manage.py migrate --noinput && \
   systemctl restart eschool-structure"
\end{lstlisting}

\section{Vérification Post-Déploiement}

\begin{lstlisting}[style=python, title=Commandes de vérification]
# Verifier le statut HTTP
curl -s -o /dev/null -w "%{http_code}\n" https://eschool-rdc.pro/

# Verifier les logs Gunicorn
journalctl -u eschool-structure --no-pager -n 20

# Verifier le service
systemctl status eschool-structure
\end{lstlisting}


% ============================================================
\chapter{Guide du Développeur}
% ============================================================

\section{Comment Ajouter un Nouveau Modèle}

\begin{enumerate}
  \item Définir la classe dans \texttt{models.py} avec \texttt{db\_table} explicite
  \item Ajouter les champs avec les contraintes (\texttt{unique\_together}, FK...)
  \item Exécuter \texttt{python manage.py makemigrations}
  \item Vérifier la migration générée, ajuster si nécessaire (cf. migration \texttt{0028})
  \item Exécuter \texttt{python manage.py migrate}
\end{enumerate}

\section{Comment Ajouter un Endpoint API}

\begin{enumerate}
  \item Créer la vue dans \texttt{views.py} avec le décorateur \texttt{@require\_http\_methods}
  \item Utiliser le pattern Upsert si applicable
  \item Toujours retourner un \texttt{JsonResponse} avec \texttt{success: True/False}
  \item Enregistrer l'URL dans \texttt{urls.py}
  \item Tester avec \texttt{curl} ou depuis le navigateur
\end{enumerate}

\section{Comment Créer une Migration Personnalisée}

Lorsque l'état de la base de production diffère du schéma Django attendu :

\begin{lstlisting}[style=python, title=Pattern RunSQL pour migrations de production]
from django.db import migrations

class Migration(migrations.Migration):
    dependencies = [('structure_app', '0027_previous')]
    
    operations = [
        migrations.RunSQL(
            sql="""
            -- Ajout conditionnel de colonne
            SET @col_exists = (
                SELECT COUNT(*) 
                FROM information_schema.COLUMNS 
                WHERE TABLE_SCHEMA=DATABASE() 
                AND TABLE_NAME='cours' 
                AND COLUMN_NAME='section_id'
            );
            SET @sql = IF(@col_exists = 0,
                'ALTER TABLE cours ADD COLUMN section_id INT NULL',
                'SELECT 1');
            PREPARE stmt FROM @sql;
            EXECUTE stmt;
            """,
            reverse_sql="ALTER TABLE cours DROP COLUMN section_id"
        ),
    ]
\end{lstlisting}


% ============================================================
% ANNEXES
% ============================================================
\appendix

\chapter{Variables d'Environnement}

\begin{center}
\begin{tabularx}{\textwidth}{l X l}
\toprule
\textbf{Variable} & \textbf{Description} & \textbf{Exemple} \\
\midrule
\texttt{DB\_HOST} & Hôte MySQL & \texttt{localhost} \\
\texttt{DB\_NAME} & Nom de la base & \texttt{countryStructure} \\
\texttt{DB\_USER} & Utilisateur MySQL & \texttt{link2ictgroup} \\
\texttt{DB\_PASSWORD} & Mot de passe MySQL & (secret) \\
\texttt{DB\_PORT} & Port MySQL & \texttt{3306} \\
\texttt{EMAIL\_HOST} & Serveur SMTP & \texttt{monecole.pro} \\
\texttt{EMAIL\_PORT} & Port SMTP & \texttt{465} \\
\texttt{EMAIL\_USE\_SSL} & Utiliser SSL & \texttt{True} \\
\texttt{EMAIL\_HOST\_USER} & Email expéditeur & \texttt{noreply@monecole.pro} \\
\texttt{SECRET\_KEY} & Clé Django & (secret) \\
\bottomrule
\end{tabularx}
\end{center}

\chapter{Liste Complète des VIEWs MySQL}

\begin{center}
\begin{tabularx}{\textwidth}{l l X}
\toprule
\textbf{VIEW} & \textbf{Type} & \textbf{Source Hub} \\
\midrule
\texttt{classes} & 1:1 & \texttt{countryStructure.classes} \\
\texttt{cours} & 1:1 & \texttt{countryStructure.cours} \\
\texttt{sessions} & 1:1 & \texttt{countryStructure.sessions} \\
\texttt{trimestres} & 1:1 & \texttt{countryStructure.trimestres} \\
\texttt{periodes} & 1:1 & \texttt{countryStructure.periodes} \\
\texttt{mentions} & 1:1 & \texttt{countryStructure.mentions} \\
\texttt{classe\_cycle} & 1:1 & \texttt{countryStructure.cycles} \\
\texttt{ecole} & Filtré & \texttt{countryStructure.etablissements} \\
\texttt{cours\_par\_classe} & Cross-JOIN & \texttt{cours} × \texttt{classe\_active} \\
\texttt{annee\_trimestre} & Universal & Multi-JOIN (5 tables) \\
\texttt{annee\_periode} & Universal & Multi-JOIN via \texttt{annee\_trimestre} \\
\bottomrule
\end{tabularx}
\end{center}

\end{document}
